\section{先把终端搬过来：SSH}
\label{sec:deploy-ssh}

便携屏能让你看到机器人上的桌面，但如果只是改参数、编译、看日志，每次搬一套外设实在有点累。安全外壳协议（Secure Shell，SSH）可以让笔记本打开机载电脑的终端：你在自己的键盘上输入，命令却在机器人上执行，结果再传回来。普通 SSH 登录只传终端数据，通常很适合远程开发。

这里先假定两台电脑已接入同一台路由器，机器人当前地址为 \texttt{192.168.50.20}，普通用户名为 \texttt{rm}。这两个值都要换成自己的；如何找到地址、用主机名代替地址以及没有路由器时怎么办，稍后在第~\ref{sec:deploy-lan}、\ref{sec:deploy-wired}~节动手。\textbf{能上网不代表两台设备之间能互通}，访客网络可能禁止客户端相互访问。

\subsubsection{装在谁上面，命令又在哪边执行}

SSH 分为客户端与服务端：笔记本运行客户端发起连接，机载电脑运行服务端接收连接。Ubuntu 上 \texttt{openssh-client} 提供前者，\texttt{openssh-server} 提供后者。Windows 10/11 通常可在 PowerShell 或终端使用 \texttt{ssh}；若找不到命令，在“可选功能”安装 OpenSSH 客户端。机载电脑装的是 Ubuntu Server 还是 Desktop，不影响使用 SSH。

\begin{example}[从笔记本进入机载电脑]
先借助便携屏，在\textbf{机器人本机终端}安装并核对服务：
\begin{tbash}
sudo apt update
sudo apt install openssh-server
sudo systemctl enable --now ssh
systemctl is-active ssh
hostname
whoami
ip -br address
sudo ssh-keygen -lf /etc/ssh/ssh_host_ed25519_key.pub
\end{tbash}
\texttt{is-active} 应输出 \texttt{active}。记下实际用户名、当前接口的 IPv4 地址，以及最后一行的 \texttt{SHA256:...} 主机指纹。若机载电脑启用了防火墙，只对调试网络放行 SSH 端口。

然后在\textbf{笔记本终端}执行，基本形式是 \texttt{ssh 用户名@地址}：
\begin{tbash}
ssh rm@192.168.50.20
hostname
whoami
pwd
exit
\end{tbash}
第一次连接会问是否信任主机密钥。把显示的指纹与刚才在机器人本机看到的指纹核对，一致再输入 \texttt{yes}，然后输入机器人上 \texttt{rm} 用户的密码；输入密码时没有星号回显是正常现象。连接后，\texttt{hostname} 应是机器人而非笔记本的名字，\texttt{pwd} 通常是远端家目录。\texttt{exit} 退出远端终端，回到笔记本。
\end{example}

如果提示 \texttt{Connection timed out}，先查地址、链路、客户端隔离与防火墙；\texttt{Connection refused} 通常表示目标拒绝该端口连接，应检查 SSH 服务和端口；\texttt{Permission denied} 则优先核对用户名、密码或密钥。用 \texttt{ssh -v rm@地址} 查看连接过程。终端提示符可自定义，容易看花眼时，主动运行 \texttt{hostname} 与 \texttt{pwd} 确认操作对象。

\begin{note}
主机密钥验证的是“你连到哪台机器”。如果以前连接过，后来出现 \texttt{REMOTE HOST IDENTIFICATION HAS CHANGED}，先核对是否重装系统、地址被另一台机器占用或连接目标错误；确认新指纹后再按文档更新记录。不要养成禁用主机密钥检查的习惯。改 SSH 配置前也应保留当前连接与本地救援入口，先运行 \texttt{sudo sshd -t} 检查配置。
\end{note}

\subsubsection{密钥和连接别名：少输一点，也更好管理}

用户密钥回答“你是谁”，与刚才用于核验服务器身份的主机密钥职责不同。私钥留在笔记本，公钥可以交给机器人；私钥密码用于保护本地私钥，也不等于机器人登录密码。下面以 Ubuntu 笔记本为例；Windows 客户端也可生成密钥，分发公钥的方法见 OpenSSH 文档。

\begin{example}[给自己的调试连接建档]
在笔记本上生成专用密钥；若该名字已存在，选新名字，不覆盖旧密钥。按提示设置私钥密码：
\begin{tbash}
ssh-keygen -t ed25519 -f ~/.ssh/id_ed25519_rm -C "rm-debug-laptop"
ssh-copy-id -i ~/.ssh/id_ed25519_rm.pub rm@192.168.50.20
ssh -i ~/.ssh/id_ed25519_rm rm@192.168.50.20
\end{tbash}
确认可登录后，退出连接，把下面一段加入\textbf{笔记本}已有的 \path{~/.ssh/config}，不要覆盖其他配置：
\begin{tbash}
Host rm-robot
    HostName 192.168.50.20
    User rm
    IdentityFile ~/.ssh/id_ed25519_rm
    IdentitiesOnly yes
    ServerAliveInterval 15
    ServerAliveCountMax 3
\end{tbash}
之后运行 \texttt{ssh rm-robot} 即可。\texttt{Host} 是本机别名，\texttt{HostName} 指定实际目标。以后可将目标改为 \path{rm-robot.local} 或 Tailscale 名称。\texttt{IdentityFile} 指明使用哪份私钥；保活参数帮助发现连接中断，不会自动重启远端程序。用 \texttt{ssh -G rm-robot} 可展开当前实际采用的客户端配置。
\end{example}

不要把私钥提交到 Git，也不要让全队共用一份个人私钥。每人使用自己的密钥，队员离队或笔记本丢失时，你就可以在机载电脑上单独撤销对应的公钥。公钥登录先在另一个终端验证成功，再考虑是否调整密码登录策略，避免把自己锁在门外。

\begin{exercise}\label{ex:deploy-ssh}
使用自己的两台设备完成一次 SSH 登录，记录客户端与服务端的主机名、用户名和主机指纹核验方式。再添加一个连接别名，以密钥登录，说明公钥、私钥、主机密钥分别保存在哪里、负责验证谁。只有一台机器时可先在虚拟机练习，并明确标注不是实机验收。
\end{exercise}

\begin{conclusion}
SSH 让你远程使用终端；它需要网络可达、远端服务启动及身份验证成功。IP 地址或可解析的名字负责找机器，主机指纹负责确认机器，用户名与用户密钥负责登录。断开 SSH 后程序是否继续运行，还取决于会话与进程的启动方式，见第~\ref{sec:deploy-ops}~节。
\end{conclusion}

安装与服务检查可查 \href{https://ubuntu.com/server/docs/how-to/security/openssh-server/}{Ubuntu OpenSSH 指南}；密钥和客户端配置参数分别见 \href{https://man.openbsd.org/ssh-keygen}{ssh-keygen 手册}、\href{https://man.openbsd.org/ssh_config}{ssh\_config 手册}。

\Needspace{10\baselineskip}
\subsubsection{用 SSH 显示远程窗口（选读）}
\label{sec:deploy-x11-forwarding}

有时只想打开机载电脑上的一个配置工具。若程序支持 \textbf{X11}，可以让它在机载电脑上运行，通过 SSH 把窗口显示在笔记本上。程序仍读取机载电脑上的文件，也仍使用那台电脑的资源。要接管机载电脑已有的整个桌面，见第~\ref{sec:deploy-desktop}~节。

X11 和 \textbf{Wayland} 都是 Linux 图形程序与显示系统之间的协议。使用 X11 时，程序通过 X 服务器显示窗口和接收输入，SSH 可以转发这条 X11 连接。使用 Wayland 时，原生程序连接到负责合成窗口和分发输入的 Wayland 合成器。\textbf{Xwayland} 是在 Wayland 会话中为 X11 程序提供兼容性的 X 服务器。

\Needspace{10\baselineskip}
\begin{example}[让机载电脑上的时钟显示在笔记本]
先确认普通 SSH 已能登录。在\textbf{机载电脑}上安装授权工具和测试程序，然后检查 SSH 服务端的配置。
\begin{tbash}
sudo apt install xauth x11-apps
sudo sshd -T | grep -E '^x11(forwarding|uselocalhost) '
\end{tbash}
预期看到 \texttt{x11forwarding yes} 和 \texttt{x11uselocalhost yes}。如果服务端尚未允许转发，在机载电脑上编辑 \path{/etc/ssh/sshd_config} 或其中引用的配置文件，启用 \texttt{X11Forwarding yes}，保留 \texttt{X11UseLocalhost yes}。保持现有 SSH 连接，先用 \texttt{sudo sshd -t} 检查配置，确认通过后在 Ubuntu 上用 \texttt{sudo systemctl reload ssh} 重载。

\Needspace{9\baselineskip}
在\textbf{笔记本的图形终端}打开一条新连接。Ubuntu 笔记本可用 \texttt{sudo apt install xauth} 安装本地授权工具。第一行在笔记本执行，其余三行在连上后的远端终端执行。
\begin{tbash}
ssh -X rm-robot
hostname
echo "$DISPLAY"
xclock
\end{tbash}
主机名应是机载电脑的名称，\texttt{DISPLAY} 通常类似 \texttt{localhost:10.0}，这个地址由 SSH 分配，保留它即可。时钟窗口应出现在笔记本上。关闭窗口后，远端终端会回到提示符。机载电脑可以只运行命令行系统，只需具备 SSH 服务、授权工具和要打开的程序。
\end{example}

没有窗口时，先按下文检查笔记本能否打开本地 X11 显示，再用 \texttt{ssh -v -X rm-robot} 检查转发请求是否被接受。远端的 \texttt{DISPLAY} 为空时，优先核对服务端配置和两端的 \texttt{xauth}。\texttt{-X} 使用受限制的转发，某些程序或本地 X 服务器可能不支持。若确认问题出在这一限制，可仅对信任的主机和程序试用 \texttt{ssh -Y rm-robot}。\texttt{-Y} 会给予远端程序访问笔记本 X11 会话的较大权限。详见 \href{https://man.openbsd.org/ssh#X11_FORWARDING}{OpenSSH 的 X11 转发说明}。

\Needspace{7\baselineskip}
\textbf{笔记本平时用 Wayland 怎么办}

先从\textbf{笔记本的桌面}打开一个新的普通用户终端。这里查的是笔记本的显示环境，所以在发起 SSH 之前执行下面的命令。
\begin{tbash}
printf 'XDG_SESSION_TYPE=%s\nDISPLAY=%s\n' "$XDG_SESSION_TYPE" "$DISPLAY"
command -v Xwayland
\end{tbash}
\texttt{XDG\_SESSION\_TYPE} 显示 \texttt{wayland} 时，\texttt{command -v Xwayland} 可检查是否找得到 Xwayland 程序，常见路径是 \path{/usr/bin/Xwayland}。有这个路径只说明找得到程序。\texttt{DISPLAY} 若有值，例如 \texttt{:0}，再实际连一次 X 显示。Ubuntu 上缺少 \texttt{xprop} 时，可在笔记本上用 \texttt{sudo apt install x11-utils} 安装。
\begin{tbash}
xprop -root >/dev/null && echo "X11 OK"
\end{tbash}

\begin{itemize}
  \item \textbf{显示 \texttt{X11 OK}。}笔记本当前终端能连接 X 服务器，可以保留当前会话，用上面的 \texttt{ssh -X} 转发。机载电脑上的程序要支持 X11。有多个图形后端时，可在远端启动命令前加 \texttt{GDK\_BACKEND=x11}（GTK）或 \texttt{QT\_QPA\_PLATFORM=xcb}（Qt），前提是已安装相应后端。
  \item \textbf{\texttt{DISPLAY} 为空，或提示 \texttt{unable to open display}。}先回到笔记本桌面，用当前登录用户打开新终端复查。若 Ubuntu 22.04 的 GNOME 会话缺少 Xwayland，在笔记本上运行 \texttt{sudo apt install xwayland}，退出桌面并重新登录，再跑一次检查。若已经安装，按所用桌面的说明确认 Xwayland 支持已启用。显示地址由桌面会话提供，手动填一个 \texttt{:0} 无法启动 Xwayland。
  \item \textbf{仍没有可用的 X11 显示。}如果程序支持原生 Wayland，可以改用下面的 \texttt{waypipe}。如果程序只支持 X11，可在笔记本切换到可用的 Xorg 会话再检查。Ubuntu 22.04 的 GNOME 登录页可在选中用户名后，通过齿轮选择 \href{https://documentation.ubuntu.com/release-notes/22.04/}{Ubuntu on Xorg}。也可按第~\ref{sec:deploy-desktop}~节连接机载电脑的远程桌面，在其支持 X11 的图形会话里启动程序。
\end{itemize}

\Needspace{6\baselineskip}
原生 Wayland 窗口可以用 \texttt{waypipe} 经 SSH 转发到笔记本的 Wayland 会话。两端都要安装 \texttt{waypipe}，并允许 SSH 转发 Unix 套接字。在 Ubuntu 22.04 可从 \texttt{universe} 软件源安装 \texttt{waypipe}。在机载电脑安装 \texttt{gtk-4-examples} 后，从笔记本的 Wayland 图形终端运行下面的命令。
\begin{tbash}
waypipe ssh rm-robot env GDK_BACKEND=wayland gtk4-demo
\end{tbash}
窗口能出现并响应操作，才算这个程序的转发成功。两端版本、图形接口及应用依赖会影响兼容性。若要看相机的实时画面或跑依赖 GPU 的图形工具，先用真实数据测延迟和帧率，再选入口。参见 \href{https://wayland.freedesktop.org/docs/book/Xwayland.html}{Wayland 的 Xwayland 说明}和 \href{https://manpages.ubuntu.com/manpages/jammy/man1/waypipe.1.html}{Ubuntu 22.04 的 waypipe 手册}。

\begin{exercise}\label{ex:deploy-x11}
记录笔记本的会话类型、本地 \texttt{DISPLAY} 和 \texttt{xprop} 的检查结果。若本地 X11 显示可用，连接 SSH 后记录远端主机名和 SSH 分配的 \texttt{DISPLAY}，打开时钟并指出程序和窗口分别在哪台电脑上。若本地检查失败，说明你会先检查哪一项，以及仅支持 X11 和支持原生 Wayland 的程序分别可走哪条路线。
\end{exercise}

\subsubsection{远程编辑：VS Code Remote -- SSH（选读）}
如果已经按卷零第 1 章在本机用过 VS Code，这里可以沿用熟悉的编辑界面，把打开的工程换成机器人上的目录。


SSH 不要求你只能用终端编辑器。第 1 章用过的 VS Code 可安装 Microsoft 的 \texttt{Remote - SSH} 扩展，在笔记本显示界面，在机载电脑读写工程、运行终端与调试器。这属于远程开发：它不会直接把机器人的整个桌面搬过来，也不会自动把所有本地文件同步过去。

\begin{example}[确认编辑与编译发生在哪台电脑]
先在普通终端验证 \texttt{ssh rm-robot} 成功，再在\textbf{笔记本}的 VS Code 扩展页安装 \texttt{Remote - SSH}。打开命令面板，选择 \texttt{Remote-SSH: Connect to Host...}，选择配置好的 \texttt{rm-robot}。首次连接会在远端安装 VS Code Server，需要可用下载通路及满足扩展要求的系统与资源；离线场地应提前准备。

连接后检查状态栏显示远程主机，再用 \texttt{File > Open Folder} 打开\textbf{远端}工程目录。在新建的集成终端运行 \texttt{hostname}、\texttt{pwd} 和 \texttt{uname -m}，应看到机器人的主机名、目录和体系结构。打开一个练习文本，加入一行并保存；在独立 SSH 终端用 \texttt{cat} 读取相同的远端路径，确认修改落在远端。C++、CMake 等开发扩展按提示安装到远端，编译器与依赖也需要在远端准备。
\end{example}

\begin{exercise}\label{ex:deploy-remote-editor}
用远程窗口打开一个练习工程，从远端集成终端构建或运行本章心跳示例。把 \texttt{hostname} 与当前路径记入日志，再指出本地窗口与远程窗口各修改哪个文件系统。若安装失败，保留扩展日志，区分 SSH 连接失败与远端 Server 安装失败。
\end{exercise}

新版本对远端 Linux 的库版本、架构与下载条件可能有不同要求，安装前核对 \href{https://code.visualstudio.com/docs/remote/ssh}{Remote -- SSH 官方教程}和其中的系统要求；这也适用于资源有限的 ARM 机载电脑。
